\documentclass[10pt,a4paper]{amsart}
\usepackage[margin=19mm]{geometry}
\usepackage{amsmath,amssymb,mathtools}
\usepackage[scr=rsfs,cal=euler]{mathalfa}
\usepackage{xfrac}
\usepackage[unicode,hidelinks,bookmarks=false]{hyperref}
\newcommand{\ie}{i.e.\ }
\newcommand{\cf}{cf.\ }
\newcommand{\resp}{respectively\ }
\newcommand{\Subsection}{\subsection}
\providecommand{\nobreakdash}{\nobreak\hbox{-}}
\setlength{\emergencystretch}{2em}
\multlinegap0pt
\theoremstyle{plain}
\newtheorem{theorem}{Theorem}
\newtheorem{definition}[theorem]{Definition}
\newcommand{\CA}{\mathcal{A}}
\newcommand{\Cab}{\mathrm{Cab}}
\newcommand{\Tw}{\mathcal{TW}}
\newcommand{\cHFL}{\mathcal{H\!F\! L}}
\newcommand{\F}{\mathbb{F}}
\newcommand{\Z}{\mathbb{Z}}
\newcommand{\Id}{\mathrm{Id}}
\newcommand{\UU}{\mathbf{U}}
\newcommand{\s}{\mathfrak{s}}
\newcommand{\tw}{\mathbf{tw}}
\newcommand{\w}{\mathbf{w}}
\def\gr{\textup{gr}}
\newcommand{\newword}[1]{\emph{\textbf{#1}}}
\title[The module structure over the cable algebra]{Colored knot Floer homology: structures and examples: the module structure over the cable algebra (Theorem 5.1)}\author{Akram Alishahi, Eugene Gorsky, Beibei Liu}
\date{Source: arXiv:2508.21776v1 (CC BY 4.0)}
\begin{document}
\begin{abstract}
Inspired by the $S^n$ colored version of Khovanov and Khovanov-Rozansky homology, we define a colored version of knot Floer homology by studying the colimit of a directed system of link Floer homology with infinite full twists. Specifically,  our $n$-colored knot Floer homology of a knot $K$ is then defined as the colimit of the link Floer homology of $(n, mn)$-cables of  $K$ by fixing $n$ and letting $m$ goes to infinity. We show that the colimit of the infinite full twists is a module over the colored knot Floer homology of the unknot.  In addition,  we give an explicit description of colored  Heegaard Floer homology for L-space knots, and maps for colored knot Floer homology of crossing changes. 


\end{abstract}
\maketitle

\noindent\emph{Source and licence.} Colored knot Floer homology: structures and examples. Version arXiv:2508.21776v1 (CC BY 4.0), distributed under the Creative Commons Attribution 4.0 International licence, http://creativecommons.org/licenses/by/4.0/, as recorded in the arXiv licence metadata for this exact version and shown on the abstract page https://arxiv.org/abs/2508.21776v1. Full rights notice: licenses/arxiv-2508.21776-CC-BY-4.0.txt.\par\smallskip

\noindent\textit{Editorial scope.} Counted unit: the module-structure theorem for cabled colored knot Floer homology, printed as Theorem 5.1 of the article (source0.tex lines 2448--2451, source label \texttt{thm: cabled module}), delivered together with the article's complete original proof of it (source0.tex lines 2453--2523, one \texttt{proof} environment of 71 lines that constructs the action, verifies the linear and the quadratic relations and closes, with no gap and no deferral). No mathematics is written, repaired or re-derived: statement and proof are the article's own text line for line, in the article's own notation ($L$, $L_m$, $\Tw_D(L)$, $\CA_n$, $a_k$, $\phi_k$, $\cHFL$, $\Cab_n(O)$, $R_{UV}$, $\F$, the Alexander, Maslov and twist gradings). Required context retained uncounted: the article's abstract, reproduced verbatim as front matter (line 191); the article's own setting for the cable algebra (lines 2108--2114), namely the direct sum $\Cab_n(O):=\Tw(O_n)=\bigoplus_{m\ge 0}\cHFL(T(n,nm))$ with its three gradings and the cobordism maps $\phi_k$ that insert a full twist; the article's own definition of the $n$-strand cable algebra $\CA_n$ with its commuting generators $a_0,\dots,a_{n-1}$ and its linear and quadratic relations, together with the Alexander, Maslov and twist degrees of the generators (lines 2117--2139); the article's sentence recording that the relations are homogeneous in the three gradings and the statement of the article's Theorem 4.6, the unlink case, which the counted proof invokes twice (lines 2141--2146); and the article's section heading and the definition of the space $\Tw_D(L)=\bigoplus_{m\ge 0}\cHFL(L_m)$ with its Alexander, Maslov and twist gradings (lines 2432--2441). Every definition, numbered display, label and structural statement of those lines is retained, unchanged. The wrapper adds one editorial section heading, ``The cable algebra and the case of the unknot'', before the first retained excerpt so that the article's own subsection heading in that excerpt does not print without a parent section; that heading is wrapper matter and is not the article's text. Omitted from this package and not redistributed: the article's front matter and titlepage block (its authors, addresses, e-mail addresses and acknowledgements of support), its keywords, the whole of its introduction with the definitions of the colored invariant, of the normalized Alexander degree and of the coloured homology, with its announced Theorems 1.1 to 1.11, its Conjecture and its Remarks (lines 199--393), its background sections on link Floer homology and on the algebra of the unknot (lines 394--2107), the parts of its sections 4 and 5 that lie between the retained excerpts (lines 2115--2116, 2140, 2147--2431, 2442--2447, 2452), the proof of its Theorem 4.6 (lines 2148--2193), which is not retained because it depends on two statements of the article that are not retained and would leave unresolved references, the remainder of its sections 4 and 5 (lines 2194--2431 and 2524--3129), its sections on applications and on crossing changes and all of its later material (lines 3130--3370), its bibliography (lines 3371--3483) other than the two entries transcribed below, and its closing \texttt{\textbackslash end\{document\}}. Nothing inside a retained excerpt is omitted or altered; the unit's own text outside the excerpts is the wrapper matter described here. In particular the article's one-line definition of $L_m$ as the link obtained by inserting $m$ full twists in $L$ (line 235, inside the introduction, which is not retained) is not carried; the retained excerpt at lines 2436--2441 recalls the construction of $\Tw_D(L)$ from the homology of the links $L_m$, and the counted proof itself describes the cobordisms from $L_m$ to $L_{m+1}$ that insert the twists, so no definition used by the counted statement or its proof is left without a referent in the unit. Citations: the retained excerpts carry six active citation commands, which cite exactly two keys of the article's own bibliography, $\mathtt{AGL}$ and $\mathtt{Zemke}$ (a seventh citation command sits inside the article's own commented-out line and stays a comment in the delivery). Those two entries are transcribed verbatim from the article's own in-source \texttt{thebibliography} (source0.tex lines 3371--3483, 41 numbered entries), and they carry the article's own printed numbers, so the delivered citation marks read as the article's do (\texttt{[1, Proposition 3.13]} and \texttt{[38]}). No other entry of the article's bibliography is transcribed or redistributed. Reference adaptation: none was needed and none was made. Every reference inside the delivered unit resolves inside it: the labels \texttt{eq: linear} and \texttt{eq: quadratic} of the two retained relation displays, which the counted proof cites five times; and the label \texttt{thm: cabled algebra} of the retained Theorem 4.6, which the counted proof cites twice. The labels \texttt{eq:CabU} and \texttt{thm: cabled module} are retained as the article's own anchors and are not cited anywhere in the unit, which produces no unresolved reference. No unresolved reference or citation is produced. Printed slips kept literal: no printed word, symbol, hypothesis or number of the counted statement or of its proof was altered. In particular the article's own author-comment macros inside the proof (the commented-out remarks in its \texttt{\textbackslash AA} and \texttt{\textbackslash EG} macros) are delivered as the article's own comments and are never expanded; the proof's own phrase ``For details, see e.g.\ \texttt{\textbackslash cite\{AGL\}}'' is commented out in the article and stays commented out; the article's own spelling of the cohomology functor as a spaced-out script symbol, its use of \texttt{\textbackslash newword}, and its own names for the decorated cobordisms $\mathcal{C}_m$, $\mathcal{C}'_m$, $\mathcal{B}_m$, $\mathcal{C}_{b,m}$ and $\mathcal{C}_{O_n,m}$ are kept. Layout and class: the article's own class is \texttt{amsart}, the same class as the wrapper, with no options given and no shipped class or style file, so no publisher class is redistributed or used. The article's preamble loads \texttt{tikz}, \texttt{tikz-cd}, \texttt{pgfplots}, \texttt{pinlabel} and \texttt{combelow} among others; no retained excerpt contains a picture, a diagram or a graphics-inclusion command, so the wrapper loads none of them and the delivered unit needs no drawing package. The wrapper declares the theorem and definition environments the retained statements need, and reproduces the article's own macro definitions that the retained text uses ($\CA$, $\Cab$, $\Tw$, \texttt{\textbackslash cHFL}, $\F$, $\Z$, $\Id$, $\UU$, \texttt{\textbackslash gr}, \texttt{\textbackslash tw}, \texttt{\textbackslash w}, \texttt{\textbackslash newword} and $\mathfrak{s}$). The article numbers its theorem-like environments inside sections and shares one counter among them; the wrapper uses the same single shared counter without the section dependence, so the retained cable-algebra definition prints as Definition 1 of the unit, the retained unlink theorem as Theorem 2 and the counted theorem as Theorem 3, where the article prints them as Definition 4.5, Theorem 4.6 and Theorem 5.1. The retained relation displays print as (2) and (3) of the unit, where the article prints them as (22) and (23). The wrapper's measure is A4 at 10pt where the article uses its class default; that is a page-appearance difference of the wrapper only, and no formula, symbol, hypothesis or argument changes. Assets: the article has sixteen figure environments, sixteen in-source TikZ pictures and thirteen distinct graphics-inclusion targets, two of which (\texttt{localcoef} and \texttt{localdisks}) have no file in the e-print at all; none of those items lies in a retained span, no retained span contains a graphics-inclusion command, a picture or a diagram, the package declares no asset dependency, and no image file is copied into or redistributed with it. All mathematics of the unit is native text with no image dependency. Licence: version arXiv:2508.21776v1 of this article is distributed under the Creative Commons Attribution 4.0 International licence, http://creativecommons.org/licenses/by/4.0/, as recorded in the arXiv licence metadata for this exact version and shown on the abstract page https://arxiv.org/abs/2508.21776v1. A full rights notice is delivered at licenses/arxiv-2508.21776-CC-BY-4.0.txt. Characters: the retained excerpts contain no character outside ASCII; the two non-ASCII characters of the article's source file (one no-break space and one en dash) both lie outside the retained excerpts, so no encoding substitution, glyph replacement or missing-character risk arises.
\section{The cable algebra and the case of the unknot}
\subsection{Cable algebra and cobordism maps} Next, we explicitly describe the cobordism maps $\phi_k$ corresponding to adding full twists  for $T(n,mn)$. In order to do that, we consider the direct sum of homology for all $n$-strand cables of the unknot with nonnegative integer slopes:
\begin{equation}\label{eq:CabU}
\Cab_n(O): =\Tw(O_n)=\bigoplus_{m=0}^{\infty}\cHFL(T(n,nm))
\end{equation}
where $O$ is the unknot and  $O_n$ is the $n$-component unlink. 

Each summand is $\Z^n\oplus \Z$ graded by Alexander multi-grading ($\Z^n$) and Maslov grading ($\Z$). Moreover, $\Cab_n(O)$ has an additional $\Z$ grading given by the cabling slope $m$, we call it \emph{twist grading} and denote it by $\tw$. Note that cobordism maps $\phi_k$, corresponding to adding a full twist, act naturally on $\Cab_n(O)$ with twist grading $1$ i.e. they map $\cHFL(T(n,mn))$ to $\cHFL(T(n,(m+1)n))$. 

\begin{definition}
The $n$-strand \newword{cable algebra} $\CA_n$ is defined as a $\Z^n\oplus \Z\oplus \Z$ graded algebra  over $R_{UV}$ with commuting generators $a_0,\ldots,a_{n-1}$ and the following relations:
\begin{itemize}
\item Linear relations
\begin{equation}
\label{eq: linear}
U_Ia_{k-1}=V_{\overline{I}}a_k
\end{equation}
where $I\subset \{1,\ldots,n\}$ is an arbitrary subset with $|I|=k$ and $\overline{I}=\{1,\ldots,n\}\setminus I.$ 
%\AA{Change $K$ to $I$, because $K$ i the knot}\EG{done}
\item Quadratic relations
\begin{equation}
\label{eq: quadratic}
a_ia_j=\UU^{k\ell-ij}a_{k}a_{\ell}
\end{equation} 
whenever $i+j=k+\ell$ and $i\le k\le \ell\le j$.
\end{itemize}
As above, the gradings correspond to the Alexander ($\Z^n$) and Maslov ($\Z$) gradings, along with an additional  $\Z$ \emph{twist grading}. The generator $a_k$ has Alexander grading 
$$
A(a_k)=\left(\frac{n-1}{2}-k,\ldots, \frac{n-1}{2}-k\right),
$$
Maslov grading $\gr_{\w}(a_k)=-k^2-k$ and twist grading $\tw(a_k)=1$.
\end{definition}

It is easy to check that the relations are homogeneous with respect to all three gradings.

\begin{theorem}
\label{thm: cabled algebra}
The homology $\Cab_n(O)$ is a $\Z^n\oplus \Z\oplus \Z$ graded module over the algebra $\CA_n$ where the generators $a_k$ act by cobordism maps $\phi_k$. Furthermore, $\Cab_n(O)$ is a free $\CA_n$-module generated by a single class $1\in \cHFL(T(n,0))$.
\end{theorem}

\section{Module structures of Colored knot Floer homology}

In this section, we use the algebra $\CA_n$ and its localization to study the module structure for cabled and colored knot Floer homology. 

\subsection{Action of cobordism maps: general case}
Recall that for an $n$-component link $L$ in $S^3$ such that every component intersects a disk $D$ positively exactly once, 
$$
\Tw_D(L)=\bigoplus_{m=0}^{\infty}\cHFL(L_m).
$$
As before, this is a $\Z^n\oplus \Z\oplus \Z$ graded vector space with Alexander, Maslov and twist gradings.

\begin{theorem}
\label{thm: cabled module}
The space $\Tw_D(L)$ is a triply graded module over the algebra $\CA_n$.
\end{theorem}

\begin{proof}
We have cobordism maps $\phi_k:\cHFL(L_m)\to \cHFL(L_{m+1})$. We construct the action of the algebra $\CA_n$ such that the generators $a_k$ act on $\Tw_D(L)$ via the maps $\phi_k$. For this to make sense, we need to verify that the maps $\phi_k$ pairwise commute and satisfy the relations \eqref{eq: linear} and \eqref{eq: quadratic}. 
We closely follow the proof of \cite[Proposition 3.13]{AGL}. 

We would like to compare the cobordism maps from $\cHFL(L_m)$ to $\cHFL(L_{m+1})$ and the similar maps from $\cHFL(T(n,mn))$ to $\cHFL(T(n,(m+1)n))$. To distinguish them, we denote them respectively by $\phi_k^{L}$ and $\phi_k^{O}$. %\AA{When we have a single $\phi$ shouldn't we have $j=1$?}

Let us denote by $L'_{m}$ the link $L_{m}$ with an extra pair of basepoints $w'_i,z'_i$ per component. We extend the coloring $\sigma$ on $L_{m}$ to $L'_{m}$ so that its codomain is $P'=P\sqcup \{p'_1, \cdots, p'_n\}$
and $\sigma'(w'_i)=p'_i, \sigma'(z'_i)=\sigma(z_i)$.
Then $\sigma'$ restricts to a coloring of $L_{m}$ with codomain $P'$, and for simplicity, we still denote it by $\sigma'$. The ring isomorphism $\mathcal{R}_P^-\cong \F[U_1, \cdots, U_n, V_1, \cdots V_n]$ extends to an isomorphism 
$$\mathcal{R}_{P'}^-\cong \F[U_1, \cdots, U_n, V_1, \cdots, V_n, U'_1, \cdots , U'_n]$$
by sending formal variables $X_{p'_i}$ to $U'_i$. Then, we have the following isomorphisms
$$\cHFL(L_m^{\sigma'})\cong \cHFL(L_m)\otimes_{\F} \F[U'_1, \cdots, U'_n] $$
and 
$$\cHFL(L_m^{'\sigma'})\cong \cHFL(L_m^{\sigma'})/ \langle U_1-U'_1, \cdots, U_n-U'_n\rangle \cong \cHFL(L_m).$$
Let $\mathcal{S}_m$ be the quasi-stabilization cobordism from $L_m$ to $L'_{m}$. Then, the composition of the induced cobordism map from $\cHFL(L_m)$ to $\cHFL(L'_m)$ with the latter isomorphism is identity.

%Under the later isomorphism, the induced map from $\cHFL(K_{n, mn})$ to itself using the quasi-stabilization cobordism from  $\mathcal{S}_m$ be the  from $K_{n, mn}$ to $K'_{n, mn}$ is identity. 


%cobordism maps $S$ corresponding to the quasi-stabilization cobordisms $\mathcal{S}_m$ from $K_{n, mn}$ to $K'_{n, mn}$,  and by this identification, the isomorphism from $\cHFL(K_{n, mn})$  to itself is the identity map. For details, see e.g. \cite{AGL}.  



Observe that $L_m$ can be obtained as band connected sum of $L$ and $T(n,mn)$ where we add one band per link component. Consider the following cobordisms: 

\begin{itemize}
\item $\mathcal{C}_{m}$ is the decorated cobordism from $L_m$ to $L_{m+1}$ obtained from the $(-1)$-surgery on an unknot bounding a disk intersecting each component of $L_m$ positively at exactly one point,
%which is a meridian of $K_{n, mn}$. 
\item $\mathcal{C}'_m$ is the induced decorated cobordism from $L'_{m}$ to $L'_{m+1}$ similar to $\mathcal{C}_m$,
\item $\mathcal{B}_m$ from $L_{m}$ to $L_{m}\sqcup O_n$ corresponding to birth of an $n$-component unlink,
\item $\mathcal{C}_{b,m}$ from $L_{m}\sqcup O_n$ to $L'_{m}$ is given by (decorated) band attachment. More generally, band attachment gives a cobordism
$\mathcal{C}_{b,m,j}$ from $L_{m}\sqcup T(n,jm)$ to $L'_{m+j}$,
\item $\mathcal{C}_{O_n,m}$  is the cobordism from $L_{m}\sqcup O_n$ to $L_{m}\sqcup T(n,n)$ given by the 2-handle attachment along the (-1)-framed unknot for inserting a full twist in $O_n$.
\end{itemize}

Define
$$\tilde{\mathcal{C}}_m=\mathcal{C}'_m\circ \mathcal{S}_m=\mathcal{S}_{m+1} \circ \mathcal{C}_m.$$
Under the aforementioned isomorphism $\cHFL(L_{m+1}^{'\sigma'})\cong \cHFL(L_{m+1})$, the homomorphism induced by the cobordism map $F_{\tilde{\mathcal{C}}_m, \s_k}$
 from $\cHFL(L_{m})$ to $\cHFL(L_{m+1})$ is equal to $\phi^{L}_k$. Note that $\mathcal{S}_m$ can be decomposed as the cobordism $\mathcal{B}_m$ containing $n$ births from $L_{m}$ to $L_{ m}\sqcup O_n$ followed by $n$ band attachments cobordism $\mathcal{C}_{b,m}$. 

By arguing as in \cite[Proposition 3.13]{AGL} (see also Proposition 3.8 in \cite{AGL}) we can isotope the attaching circle of the 2-handle in $\tilde{\mathcal{C}}_m$ so that we can
change the order of 2-handle attachment and band attachments, which shows that the composition
$$
\tilde{\mathcal{C}}_m=\mathcal{C}'_{m}\circ \mathcal{C}_{b,m}\circ \mathcal{B}_{m}
$$
is isotopic to the composition
$$
\mathcal{C}_{b,m,1}\circ \mathcal{C}_{O_n,m}\circ \mathcal{B}_{m}.
$$

Therefore by naturality (\cite{Zemke}) we get
$$
\phi_k^{L}=F_{\tilde{\mathcal{C}}_m, \s_k}=F_{\mathcal{C}_{b,m,1}}\circ F_{\mathcal{C}_{O_n,m},\mathfrak{s}_k}\circ F_{\mathcal{B}_{m}}=F_{\mathcal{C}_{b,m,1}}\circ \left(\Id\otimes \phi_k^{O}\right)\circ F_{\mathcal{B}_{m}}.
$$
Here, $\phi_k^{O}: \cHFL(O_{n})\rightarrow \cHFL(T(n,n))$ denotes the full twist cobordism map. Since all the maps are $R_{UV}$-linear and  $\phi_k^{O}$ satisfy linear relations \eqref{eq: linear} by Theorem \ref{thm: cabled algebra}, the maps $\phi_k^{L}$ satisfy \eqref{eq: linear} as well. 

Next, we need to study the compositions of two maps $\phi_{\ell}^{L}: \cHFL(L_{m})\rightarrow\cHFL(L_{m+1})$ and $\phi_{k}^{L}: \cHFL(L_{ m+1})\rightarrow\cHFL(L_{ m+2})$
$$
\phi_{k}^{L}\circ \phi_{\ell}^{L}=
\phi_{k}^{L}\circ
F_{\mathcal{C}_{b,m,1}}\circ \left(\Id\otimes \phi_k^{O}\right)\circ F_{\mathcal{B}_{m}}.
$$



By isotopying the attaching circle of the second $2$-handle, changing the order of band attachments and the second 2-handle attachment, and using the naturality of link Floer homology we have
$$
\phi_{k}^{L}\circ \phi_{\ell}^{L}=F_{\mathcal{C}_{b,m,2}}\circ \left(\Id\otimes \left(\phi_{k}^{O}\circ\phi_{\ell}^{O}\right)\right)\circ F_{\mathcal{B}_{m}}.
$$
Consequently, since \eqref{eq: quadratic} and commutation relations hold for the unlink by Theorem \ref{thm: cabled algebra}, they hold for $L$ as well. 
\end{proof}

\begin{thebibliography}{99}
\bibitem[1]{AGL} A. Alishahi, E. Gorsky, B. Liu. Splitting maps in link Floer homology and integer points in permutahedra. arXiv:2307.07741
\bibitem[38]{Zemke} I. Zemke. Link cobordisms and functoriality in link Floer homology. J. Topol. 12 (2019), no. 1, 94--220.
\end{thebibliography}
\end{document}
